% Modelo de ficha de leitura. Cada ficha é um \chapter com blocos fixos.
% Regras: vírgula no lugar de travessão; ortografia atual fora de citação;
% citações diretas curtas, sempre com página; nada de "cabe ressaltar".

\newcommand{\ficha}[2]{\chapter[#1]{#1}\label{ficha:#2}\nocite{#2}}

% Bloco de identificação: referência ABNT completa e dados de leitura.
\newenvironment{identificacao}{%
  \section*{Referência}\begin{small}}{\end{small}}

\newcommand{\bloco}[1]{\section*{#1}}

% Citação-chave com página: \chave{texto}{p. 12}
\newcommand{\chave}[2]{%
  \begin{quote}\small ``#1'' (#2)\end{quote}}

% Ponto de diálogo com as fontes primárias: \dialogo{o que procurar nos jornais}
\newenvironment{dialogo}{%
  \section*{Pontes para as fontes primárias}\begin{itemize}[leftmargin=*,nosep]}{\end{itemize}}

% Uso no capítulo: em qual parte da estrutura entra.
\newcommand{\usonocapitulo}[1]{\section*{Uso no capítulo}#1}
